\section{Discussion}
\subsection{What P001 establishes}
P001 establishes that the repository can freeze a deletion request, construct an exact
retained complement, prevent deleted observations from entering exact retraining, execute
multiple independent replicas, store raw predictions and checkpoints with hashes, derive
reference statistics, and generate manuscript values from those artifacts. It also shows
that stochastic variation is empirically visible even in a small deterministic pilot.
These are infrastructure and measurement findings, not evidence that an approximate
method satisfies a general definition of unlearning.

\subsection{Identifiability before equivalence}
The pilot exposes an identifiability problem rather than a method ranking. Ordinary IID
deletions may leave the trained function almost unchanged relative to ordinary retraining
variation. In such a regime, approximate/exact proximity is compatible both with genuine
removal and with doing essentially nothing. Stronger class-conditioned, influential, and
controlled-memorization regimes are therefore required. This interpretation remains
restricted to assay design and does not establish the behavior of any unlearning method.

This ordering matters for equivalence testing. An equivalence interval contained within a
margin can support similarity only for the measured endpoint and only if the endpoint is
capable of detecting the relevant intervention. Otherwise scientifically different
procedures may satisfy the same surface criterion. C001 therefore requires both a
reference-sensitivity gate and prospectively frozen equivalence margins.

The class-conditioned calibration makes this distinction concrete. JS divergence detected
the original-versus-full-deletion intervention for two classes, yet the exact/exact 95th
percentile did not detect the prospectively frozen one-percent residual-class violation.
Identifiability of a large intervention is therefore not sufficient to validate a tight
equivalence margin. The failed margin is retained as a negative design result rather than
widened after observing the calibration.

\subsection{Why multiple exact replicas matter}
Across the 15 exact seeds, test accuracy varies by more than four percentage points and
forget accuracy by ten percentage points. Selecting one replica can therefore change the
apparent distance or utility cost assigned to an approximate method. An ensemble exposes
this dependence and supports empirical percentiles, bootstrap intervals, or equivalence
tests at the model-seed level. It does not make the finite ensemble identical to the full
training distribution, but it makes uncertainty visible.

\subsection{Auditability as a methodological property}
Output similarity alone cannot reveal which data were used or whether the declared
algorithm was actually executed. The repository therefore couples behavioral evidence
with procedural evidence: immutable manifests, exclusion assertions, configuration and
artifact hashes, append-only raw records, failure metadata, and claim-to-artifact links.
This does not solve every auditing problem, but it narrows the gap between a reported result
and the executable process that produced it.

\subsection{Efficiency conditional on fidelity}
The aggregate approximate-method runtime is smaller than aggregate exact retraining in
P001, but the comparison is a small hardware-specific observation. Computational savings
become scientifically meaningful only after an assay can distinguish the relevant
counterfactual and a method's fidelity is characterized. The planned compute--fidelity
frontier therefore reports cost jointly with multiple fidelity dimensions rather than as a
standalone speedup claim.

\subsection{Implications for C001}
The next study should retain class-conditioned deletion as a distinguishability stress
regime while treating equivalence as unresolved. Influence-selected observations,
controlled memorization, and corrupted-label removal remain candidate regimes. Multiple
original seeds are needed to estimate original/exact variation rather than conditioning on
one original model. Metric selection, uncertainty reporting, multiplicity control,
exclusions, and stopping rules must be registered before confirmatory training begins.
